This is within \eqref{eq:v5-global-gowers}. Let $f'_{j,b}$ be the
model supplied by Corollary~3 at cutoff $n^{\rho_*}$. Its bound is
$0\leq f'_{j,b}\leq B_*-1$, where $B_*\geq2$ depends on the fixed
ray field, complexities and on $\rho_*$, but not on $v,n,b$ or the polynomials
within the stated coefficient bounds. This is the bounded model of
Tao--Ziegler's Proposition~4: its construction first rescales the
prime weight by a fixed positive constant and then divides the model
by that constant. In particular, it need not be bounded by one.
Since $0\leq\mu_j(b)\leq1$, the function
\[
 F_{j,b}=\frac{f'_{j,b}-\mu_j(b)}{B_*}
\]
is bounded in magnitude by one. The triangle inequality, the global
part of Corollary~3, and \eqref{eq:v5-global-gowers} show that its
required global norms tend to zero under the diagonal limit below.
Apply Theorem~9 to $F_{j,b}$
with parameters $d=d_0=2$, $M=M_1$, $A=\mathcal A$, and global
order $a$. Its smallness threshold is independent of $\mathcal A$;
$B_*$ has already been fixed before $v$ is chosen. Multiplication
by $B_*$ and the local part of Corollary~3 therefore show that the
local norm of $g_{j,b}$ tends to zero. Theorem~13, with the original
parameter $A$ and coarse scale $M_0$, makes every polynomial average
containing a centered factor $o(1)$. These are the reasons for choosing
the complexities before $v$.
